% Introductory and closing texts, from john-dowland-letter-modern-english.docx.

\newcommand{\AboutText}{This version keeps close to Dowland’s meaning and sequence while replacing sixteenth-century spelling, grammar and idiom with current English. Where the original permits more than one reading, the translation chooses the most likely sense and records the uncertainty in the numbered notes. The notes also explain people, institutions and expressions that are easy to misread today.}

\newcommand{\DoingText}{The letter is simultaneously a confession, a defence, an intelligence report and an indirect request for renewed trust and patronage. Dowland constructs a careful story in which youthful persuasion led him toward Catholicism; his religious reputation blocked advancement in England; foreign princes valued him; Catholic agents then offered him money, honour and help for his family; and the discovery of anti-Elizabethan activity brought him back to open loyalty. His information may be sincere, but its presentation is strategic: it explains his past, distances him from treason and asks Cecil to regard him again as a useful English subject.}

\newcommand{\SourceText}{The translation follows David Pinto’s diplomatic and expanded transcription of the letter in the critical hypertext edition hosted by the University of Birmingham. The manuscript is held at Hatfield House in the Cecil Papers, volume 172, numbers 91-93.}
